\documentclass[10pt,a4paper]{amsart}
\usepackage[margin=19mm]{geometry}
\usepackage{amsmath,amssymb,mathtools}
\usepackage[scr=rsfs,cal=euler]{mathalfa}
\usepackage{xfrac}
\usepackage[unicode,hidelinks,bookmarks=false]{hyperref}
\newcommand{\ie}{i.e.\ }
\newcommand{\cf}{cf.\ }
\newcommand{\resp}{respectively\ }
\newcommand{\Subsection}{\subsection}
\providecommand{\nobreakdash}{\nobreak\hbox{-}}
\setlength{\emergencystretch}{2em}
\multlinegap0pt
\usepackage{dsfont}
\usepackage{amssymb}
\usepackage{mathtools}
\usepackage{natbib}
\usepackage{tikz}
\usepackage{siunitx}
\usepackage{booktabs}
\usepackage{multirow}
\newtheorem{corollary}{Corollary}
\newtheorem{theorem}{Theorem}
\title{Simultaneous symplectic spectral decomposition of positive semidefinite matrices}
\author{Rudra R. Kamat, Hemant K. Mishra}
\date{Source: arXiv:2412.01492v3 (CC BY 4.0)}
\begin{document}
\maketitle

\noindent\emph{Source and licence.} Simultaneous symplectic spectral decomposition of positive semidefinite matrices. Version arXiv:2412.01492v3 (CC BY 4.0), distributed by arXiv under the Creative Commons Attribution 4.0 International licence, http://creativecommons.org/licenses/by/4.0/, as recorded in the arXiv licence metadata for this exact version and shown on the abstract page https://arxiv.org/abs/2412.01492v3. Full licence notice: \texttt{licenses/arxiv-2412.01492-CC-BY-4.0.txt}.\par\smallskip

\noindent\textit{Editorial scope.} Counted unit: the article's theorem thm:symplectic-commute-powers (source0.tex lines 413--416). The article's complete original proof of it is delivered with it (source0.tex lines 417--427, one \texttt{proof} environment of 11 lines). No mathematics is written, repaired or re-derived: the statement and the proof are the article's own lines, in the article's own notation, and no retained line is substituted or re-flowed.  Retained uncounted as the source-local context the proof needs: lines 378--380.  Nothing was omitted from any retained span, so the package carries no omission record.  No retained line contains a citation, so the wrapper carries no bibliography and no bibliography entry.  No retained line contains a figure environment, an image-inclusion command or drawing code, so no image is delivered and the package declares no asset dependency.  Editorial formatting only, in addition: an AMS \texttt{amsart} preamble that restates the article's own declarations and macros used by the retained lines, namely the environments \texttt{corollary}, \texttt{theorem} on separate counters, together with the standard packages those lines need.  The retained environments print in this document as Corollary 1, Theorem 1 in order of appearance, whereas the article prints the same items with the numbering of its own full document.  The complete native source of the article as distributed in the arXiv e-print for this exact version is bundled unchanged as \texttt{sources/169462/source0.tex}.
    \begin{corollary}\label{cor:orthosymplectic-diagonalization-williamson}
        Orthosymplectic spectral decomposition of a $2n \times 2n$ real positive semidefinite matrix $A$ with symplectic kernel exists if and only if $JA=AJ$.
    \end{corollary}

    \begin{theorem}\label{thm:symplectic-commute-powers}
        Let $A, B$ be $2n \times 2n$ real positive definite matrices.
        If $AJB=BJA$ and $AB=BA$, then we have $A^s J B^s= B^s J A^s$ for all $s \in \mathds{R}$.
    \end{theorem}

    \begin{proof}
        The condition $AJB=BJA$ implies that
        \begin{align}\label{eq:symplectic-commutativity-modified}
            B^{-1}AJ=JAB^{-1}.
        \end{align}
        Also, since $A$ and $B$ commute, the matrix $AB^{-1}$ is a symmetric positive definite matrix. 
        Therefore, combining \eqref{eq:symplectic-commutativity-modified} and Corollary~\ref{cor:orthosymplectic-diagonalization-williamson}, we get that $AB^{-1}$ is orthosymplectically diagonalizable in the sense of Williamson's theorem.
        This implies that for any $s \in \mathds{R}$ the matrix $A^sB^{-s}$ is also orthosymplectically diagonalizable in the sense of Williamson's theorem.
        By invoking Corollary~\ref{cor:orthosymplectic-diagonalization-williamson} again, we thus have $JA^sB^{-s}=A^sB^{-s}J$.
        Using the commutativity of $A$ and $B$, this simplifies to $A^s J B^s = B^s J A^s$.
    \end{proof}

\end{document}
